\documentclass[11pt]{article}
\usepackage{setspace}
\usepackage{geometry}
\geometry{margin=1in}
\setstretch{1.1}

\begin{document}

\begin{flushright}
\textbf{Zhongli Wang} \\
Phone: +49 160 9541 6015 \\
Email: nkwangzhongli@gmail.com \\
Website: https://sites.google.com/view/zhongli-wang \\
Bielefeld, Germany \\
\end{flushright}

\vspace{0.25cm}

\noindent
\textbf{Selection Committee} \\
Chair of Environmental Economics \\
Department of Business Studies and Economics \\
RPTU University Kaiserslautern-Landau \\
Kaiserslautern, Germany \\

\vspace{0.25cm}

\noindent
\textbf{Application for Postdoctoral Researcher in Environmental Economics / Economic Theory}\\

\vspace{0.15cm}

\noindent Dear Prof. Heyen and Members of the Selection Committee,\\[0.1cm]

I am writing to apply for the position of Postdoctoral Researcher in Environmental Economics / Economic Theory at RPTU University Kaiserslautern-Landau. I am close to completing my Ph.D.\ in Economics at Bielefeld University and Universit\`a Cattolica del Sacro Cuore within the Marie Sk{\l}odowska--Curie EPOC program. My research focuses on dynamic economic systems, strategic interaction and agent-based modeling. I am particularly interested in how decision-makers respond to uncertainty, information, and policy interventions in complex environments.

My doctoral dissertation develops computational models of market dynamics under regulation, consumer heterogeneity, network effects, congestion, and strategic firm investment. While the empirical context of my dissertation is the ICT industry, its methodological core is closely aligned with the advertised DFG project: decision-making under uncertainty, dynamic adaptation, information diffusion, strategic interaction, and system-level feedback. These are also central issues in socio-ecological systems that may exhibit tipping points, where early warning signals can change beliefs, incentives, and behavior before irreversible transitions occur.

A central strength I would bring to the project is my experience in modeling dynamic decision-making and learning in complex environments. In my job market paper, I study how firms make capacity-investment decisions when demand is shaped by network effects and congestion. Firms update their prices and investment levels in response to observed past outcomes, including profitability, market share, and consumer behavior. This learning-from-feedback structure is closely related to the project's focus on how early warning systems affect beliefs and decisions under tipping risks. I understand early warning systems not only as statistical indicators of approaching thresholds, but also as information structures that change agents' expectations and behavior. In this respect, I was particularly interested in Prof. Heyen's work showing that warning systems can help decision-makers update beliefs about the location of a tipping point, while the absence of a warning may also induce more risky behavior. My own modeling experience would allow me to contribute to analyzing how agents learn from observed signals, how they adapt when signals are noisy or absent, and how such learning affects system-level outcomes.

My background also fits the project's focus on one-shot and dynamic decision problems. In my dissertation, agents repeatedly update decisions in response to observed market signals rather than solving a fully static problem. This provides a useful foundation for developing models in which decision-makers decide whether and how to respond to early warning information under uncertainty. I would be especially interested in analyzing how the value of warning systems depends on the timing, reliability, and interpretation of signals, and how the behavioral effect of receiving no warning may differ from the effect of receiving an explicit warning. This distinction is important because the absence of a warning can itself become a decision-relevant signal, potentially encouraging either continued precaution or more risky behavior. 
I also bring strong numerical, programming, and broader research skills. My third dissertation paper focuses on parameter estimation in an agent-based computational model using financial and operational data from AT\&T and T-Mobile. I apply the Simulated Method of Moments and neural network-based surrogate modeling, supported by large-scale simulation pipelines implemented in Python and high-performance computing resources. Beyond my dissertation, my co-authored paper on \emph{Polytope Fraud Theory} applies supervised machine learning to detect structured anomalies in financial statements, while my work on \emph{Agentic Workflows for Economic Research} develops a framework for using LLM agents to support reproducible research pipelines. Together, these projects demonstrate my ability to formulate research questions, build computational tools, work with complex data and models, and translate technical results into publishable academic contributions. These experiences prepare me to contribute to the project's analytical and numerical modeling tasks, including solving dynamic problems, exploring comparative statics, and generating predictions that can be tested experimentally.

I am strongly motivated by the project's integration of theory and experiments. I have designed models with explicit empirical and policy interpretation, and I would be very interested in collaborating with the experimental team at the University of Augsburg. I see an important role for theory in clarifying the mechanisms through which early warning systems affect behavior and in generating testable hypotheses for experimental designs. For example, theoretical analysis could help identify when warning signals improve welfare, when they are ignored, when they trigger overreaction, and when strategic interaction changes their effectiveness.

RPTU University Kaiserslautern-Landau and the Chair of Environmental Economics provide an excellent environment for this research. I am particularly drawn to the opportunity to work under the supervision of Prof. Daniel Heyen and to contribute to a collaborative project linking economic theory and experiments. I would bring to the team a strong foundation in dynamic modeling, game-theoretic reasoning, programming, together with a collaborative attitude and a strong willingness to engage with environmental economics and experimental research.

Thank you so much for considering my application! I would be delighted to discuss how my background in economic theory, computational modeling, and decision-making in complex systems can contribute to the project.

\vspace{0.35cm}

\noindent Sincerely,\\[0.95cm]
\noindent Zhongli Wang

\end{document}
